\hrulefill

\section{Introduction}

Welcome to this data labelling guide for historical corpora annotation. This document will provide you with all the information and context necessary for your work as an annotator. This document is available to you throughout your assignment and you should refer to it if you have any questions.

We want to produce annotations to train a model to detect different types of content in historical documents. To train the model, the annotations must be as consistent as possible.

\subsection{Context}

Our goal is to create an automated system for extracting visual elements using machine learning technology. For this, we have high-quality data, annotated and manually corrected by human hands.

We are currently working on a study of the dissemination of texts by studying images. In this project, we need only interested in the visual elements and we will completely disregard text.

The page of a historical document is composed of text, illustrations, ornaments and other elements added after the creation of the object, such as glosses, comments, stamps, drawings, traces of wear.

\subsection{Task Definition}

You are an annotator and your goal is to annotate the visual elements of the page. To do this, you will draw straight rectangles around the elements that have a label. The 5 labelling tags are the following : 

\begin{itemize}
    \item 0: "Illustration"
    \item 1: "Ornament"
    \item 2: "Initial"
    \item 3: "Stamp"
    \item 4: "Table"
\end{itemize}
